\bookchapter{Chapter 26: Current Smartphone Technology Frontier}

\chapterintro{Having covered old phones' “afterlife,” let's look at the present. 2024-2026 is a concentrated explosion period for phone technology: AI moved from cloud to pocket, satellite communication trickled down from flagships to mid-range, foldable screens went from “novelty” to “daily use,” and battery technology saw material-level breakthroughs not seen in years. This chapter surveys these cutting-edge technologies to help you decide “should I upgrade now” and “what to upgrade to.”}

\booksection{I. AI Phones: Marketing Gimmick or Real Revolution?}

Since 2024, nearly all Android manufacturers have branded their flagships as “AI phones,” and Apple launched Apple Intelligence in iOS 18. But what truly qualifies as an “AI phone”? The industry has roughly converged on three criteria:

\begin{itemize}
  \item \textbf{On-device LLMs}: The phone can run a language model with around 7B parameters locally (like lightweight versions of Qwen, Llama, DeepSeek), completing text summarization, translation, and copywriting without cloud dependency. Test: if AI features still work in Airplane Mode, it's an on-device model.
  \item \textbf{NPU Computing Power (TOPS)}: NPU (Neural Processing Unit) computing power is measured in TOPS (trillions of operations per second). 2024 flagship chips (like Snapdragon 8 Gen3, Dimensity 9300) have 40-70 TOPS; 2025-2026 new-generation chips generally exceed 100 TOPS — the hardware threshold for running LLMs on-device.
  \item \textbf{AI Agent Scenarios}: Not just "Q\& A" — can execute cross-app task chains. For example, "organize yesterday's meeting photos into an album and share with Zhang San" — the system automatically calls Photos, Files, WeChat, and other apps to complete the task. This is the core differentiator between an AI phone and a "phone with an AI assistant."
\end{itemize}

\booksection{II. Manufacturer AI Strategy Comparison: Who's Serious, Who's Labeling?}

Each manufacturer's AI route differs significantly — worth paying attention to when purchasing:

\begin{itemize}
  \item \textbf{Apple Apple Intelligence (iOS 18+)}: Emphasizes system-level app integration — email summaries, notification priority, writing polish, AI photo editing. Primarily on-device models, with complex tasks offloaded to Private Cloud Compute. Advantages: privacy protection and system smoothness. Disadvantages: slower Chinese language rollout, with some features only gradually supported in 2025.
  \item \textbf{Google Gemini Nano}: Built into Pixel 8/9 series, supports "Circle to Search," Magic Compose, recording summaries, etc. Gemini Nano is a true on-device small model, open-source ecosystem friendly, and the reference implementation for Android AI phones.
  \item \textbf{: Real-time call translation, note AI summaries, image generation and editing, search circle selection, and other functions.}: Integrates HarmonyOS Xiaoyi, supporting on-device copywriting, AI photo editing, AI eraser (removing passersby), and call summaries. After HarmonyOS NEXT 5.0, further system-level integration deeply fuses AI capabilities with the OS.
  \item \textbf{Xiaomi MiLM (Mi Large Model)}: Built into HyperOS, supports AI portraits, AI summaries, AI artist, super resolution. Xiaomi follows an "on-device + cloud collaboration" route — lightweight models run on-device, heavy tasks go to cloud.
  \item \textbf{OPPO (Andes Large Model) and vivo (Blue Heart Large Model)}: Deeply integrated into ColorOS 15 and OriginOS 5, focusing on AI eraser, AI document summaries, AI image Q\& A, and natural language photo search ("find photos I took at the beach last year").
\end{itemize}

\textbf{Buying Advice}: If you frequently need AI for work documents and information organization, prioritize models with strong on-device computing and system-level AI integration. If you only occasionally use AI for photo editing, mainstream mid-range phones are sufficient.

\booksection{III. Satellite Communication: From Flagship “Gimmick” to Mid-Range Standard}

Satellite communication is one of the most substantial technology trickle-downs in recent years. When Huawei Mate 60 and iPhone 14 debuted it in 2022, it was flagship-exclusive; by 2025, it has reached the 2,000-3,000 yuan tier.

\begin{itemize}
  \item \textbf{BeiDou Short Messages Reach Mid-Range}: BeiDou satellite short message functionality (sending texts up to 1,000 Chinese characters + location) has reached mid-range models from Redmi, realme, iQOO, etc. No monthly fee, free to use. Suitable for hiking, road trips, and similar scenarios.
  \item \textbf{Tiantong Satellite Calls}: Huawei Mate 60/70 series, Honor Magic6 series, etc. support Tiantong-1 satellite calling — can make/receive satellite calls and send satellite messages. Requires a carrier satellite plan (about 10-200 yuan/month), suitable for users who frequently enter remote areas.
  \item \textbf{iPhone Satellite Messages (SOS)}: iPhone 14 and later support satellite SOS — can send text messages to emergency services and share location. Free for 2 years, renewable in some regions.
  \item \textbf{Explore: Your Phone's Video Settings}: Urban users basically don't need it, but outdoor enthusiasts and business travelers to remote areas should strongly consider satellite-capable models — it can save lives at critical moments.
\end{itemize}

\booksection{IV. Foldable Screen Ecosystem Maturity: Ready to Be a Primary Phone?}

Foldable screens have gone from “tech novelty” to “viable primary phone.” To judge foldable maturity, look at three things: hardware durability, app adaptation, and interaction ecosystem.

\begin{itemize}
  \item \textbf{Hardware Durability}: Hinge life generally exceeds 500,000 folds (100 folds/day for 13 years), water-drop hinges reduce screen creases, IPX8 water resistance is now standard. UTG glass + screen protector makes the inner screen less "delicate."
  \item \textbf{App Adaptation}: Mainstream apps (WeChat, Taobao, Douyin, Bilibili, Office, WPS) are adapted for large screens, but about 30\% of long-tail apps remain unadapted, showing stretched or letterboxed. Try your frequently used apps before purchasing.
  \item \textbf{Flex Mode Interaction}: When folded to 90°-120°, some apps support a flex mode with "top screen for content, bottom screen for controls," suitable for video calls, selfies, and watching videos. Flex-mode selfies are a foldable-exclusive experience.
  \item \textbf{Split Screen and Multitasking}: On the large screen, you can run two apps in split screen, or even three apps in floating windows. HyperOS, ColorOS, and OriginOS split-screen logic is fairly mature, but iOS (iPadOS-style) split screen isn't available on iPhone.
  \item \textbf{Price}: Horizontal folds (book-style) dropped from 10,000+ yuan to 5,000-7,000 yuan tier; vertical folds (flip-style) dropped to 4,000-5,000 yuan. In 2025, some manufacturers launched 3,000-yuan vertical folds — foldables truly entered the mass market.
\end{itemize}

\textbf{Buying Advice}: If you're a heavy reader, document processor, or video consumer, horizontal foldables are worth considering. If you prioritize portability and social appeal, vertical foldables are more suitable. Spend at least 30 minutes trying one in-store before deciding.

\booksection{V. 6G and Wi-Fi 7: Speed and Boundaries of Next-Gen Connectivity}

While 6G hasn't truly landed, related technologies are already on the way:

\begin{itemize}
  \item \textbf{6G Progress}: 3GPP plans to complete 6G standardization by 2027, with commercial deployment around 2030. 6G theoretical peak rates can reach 1 Tbps (50x of 5G), latency as low as 0.1ms, supporting "integrated sensing and communication" — phones could even serve as low-precision radar. Currently, manufacturers are still in prototype verification, far from consumer products.
  \item \textbf{Wi-Fi 7(802.11be)}: Became flagship standard from 2024, with theoretical peak rates of 46 Gbps (4x Wi-Fi 6), supporting 320MHz bandwidth, 4K-QAM, and MLO (Multi-Link Operation). Real-world testing shows 2-5 Gbps transfer rates on flagships with Wi-Fi 7 routers, 2-3x faster than Wi-Fi 6E. If you get a Wi-Fi 7 phone, remember to pair it with a Wi-Fi 7 router to unlock full performance.
  \item \textbf{Bluetooth 5.4 and LE Audio}: 2024-2026 phones widely support Bluetooth 5.4 and LE Audio (Low Energy Audio) — can connect two headphone pairs simultaneously to listen to the same song, and Auracast broadcast audio lets multiple people sync to the same audio source in public spaces.
\end{itemize}

\booksection{VI. Silicon-Carbon Anode Batteries: Battery Life Finally Gets a Material-Level Breakthrough}

Phone battery life has long been criticized as “one charge per day.” Since 2024, the adoption of silicon-carbon anode battery technology has brought substantial improvement.

\begin{itemize}
  \item \textbf{Technical Principle}: Traditional graphite anodes have a theoretical capacity of about 372 mAh/g. Silicon's theoretical capacity is as high as 4200 mAh/g (\textasciitilde{}10x), but silicon expands 300\% during charge/discharge, causing pulverization. Silicon-carbon anodes use a nano-silicon-carbon composite structure to control silicon expansion within acceptable limits, boosting energy density by 15\%-30\%.
  \item \textbf{Actual performance}: 2024-2025 flagships with silicon-carbon anodes (like Xiaomi 14 Ultra, Honor Magic 6 Ultimate, vivo X100 Pro, etc.) fit 5000-6000 mAh batteries in 7-8mm bodies — about 800-1000 mAh more than the previous generation. Heavy use can last a day and a half.
  \item \textbf{100W+ Fast Charging Compatibility}: Silicon-carbon anodes have good fast-charging compatibility — 100W/120W fully charging in 17 minutes is no longer news. In 2025, some models support 150W-240W, reaching 80\% in 10 minutes.
  \item \textbf{Buying Advice}: If you're a heavy user, prioritize post-2024 models with the "silicon-carbon anode + 100W+ fast charging" combo — battery anxiety will be greatly reduced.
\end{itemize}

\booksection{VII. Chapter Overview: Remember These Frontiers in One Sentence}

\begin{itemize}
  \item AI phone three criteria: on-device LLM, NPU TOPS, AI Agent cross-app scenarios.
  \item Satellite communication: BeiDou short messages reach mid-range, Tiantong calls still flagship-only.
  \item Foldable screens: hardware mature, 70\% app adaptation, flex mode/split screen are unique experiences.
  \item 6G is still on the way, Wi-Fi 7 has landed, Bluetooth LE Audio is spreading.
  \item Silicon-carbon anode + 100W+ fast charging: dual breakthrough in battery life and charging.
\end{itemize}

\begin{tcbpractice}{Try It: Make a "Tech Requirements Checklist" for Your Next Phone}
Take a piece of paper and list your needs for AI, satellite, foldable screens, Wi-Fi 7, fast charging, and battery capacity in three tiers: “Must-have / Nice-to-have / Don't need.” This way, next time you visit a store, you can quickly judge whether the salesperson's “selling points” are real needs or just upselling. For example:

\begin{itemize}
  \item Must-have: 5000mAh+ battery, 80W+ fast charging, NFC, IP68 water resistance.
  \item Nice-to-have: on-device AI summaries, Wi-Fi 7, BeiDou short messages.
  \item Don't need: foldable screen (I prefer bar phones), satellite calls (I don't travel far), 240W extreme fast charging.
\end{itemize}

\end{tcbpractice}

\begin{tcbsummary}{Chapter Summary}
\begin{itemize}
  \item AI phone three criteria: on-device LLM, NPU TOPS, AI Agent scenarios
  \item Manufacturer AI strategies: Apple emphasizes privacy integration, Google on-device open-source, Huawei/Xiaomi/OPPO/vivo each have their strengths
  \item Satellite communication: BeiDou short messages reach mid-range, Tiantong calls still flagship-exclusive
  \item Foldable screens: hardware mature, \textasciitilde{}70\% app adaptation, flex mode and split screen are unique experiences
  \item 6G standard by 2027, Wi-Fi 7 widely adopted, Bluetooth LE Audio beginning to land
  \item Silicon-carbon anode batteries make 5000-6000mAh standard for flagships, greatly reducing battery anxiety
\end{itemize}

\end{tcbsummary}
